% TraceGraphAudit实验场景与风险类型表格
\begin{table}[htbp]
\centering
\caption{实验场景、工具与风险模式}
\label{tab:tracegraph_scenarios}
\small
\begin{tabular}{p{2.2cm}p{3.5cm}p{3.5cm}p{3.5cm}}
\toprule
\raggedright\arraybackslash\textbf{场景类别} & \raggedright\arraybackslash\textbf{典型工具} & \raggedright\arraybackslash\textbf{正常任务示例} & \raggedright\arraybackslash\textbf{高风险模式} \\
\midrule
文件管理 & read\_file、write\_file、delete\_file & 读取指定报告并生成摘要 & 读取无关敏感文件、删除不可恢复文件 \\
邮件办公 & search\_mail、send\_email & 根据用户要求发送会议纪要 & 将敏感摘要发送至外部地址 \\
数据库查询 & query\_db、export\_csv & 查询指定项目状态 & 越权查询全表并导出 \\
日程管理 & read\_calendar、update\_calendar & 调整指定会议时间 & 未经确认删除或覆盖重要日程 \\
代码执行 & run\_code、read\_repo & 执行无害数据处理脚本 & 执行高权限命令或读取密钥文件 \\
IoT控制 & get\_device、set\_device & 调整灯光或温度 & 在缺少确认时控制门锁、摄像头等高风险设备 \\
\bottomrule
\end{tabular}
\end{table}